\documentclass[10pt]{article}
\usepackage{icmlko}

\icmlmeta{IP 파운데이션 월드모델}{82}{빼면 오른다는 말은 산술이었다}{애니를 빼면 판정치가 오르는데 펀딩을 빼면 더 오른다 --- 판정치로는 도메인을 못 고른다}

\begin{document}

\twocolumn[
\icmlheader

\begin{abstract}
노트 81이 숙제를 남겼다 --- 애니(라프텔)는 팝업 출처로 $+$0.124(꼴찌)이고
세계애니가 같은 작품을 $+$0.474로 맞히니 라프텔을 뺄지 정하라. 씨앗 넷 규약으로
쟀더니 \textbf{4/4 채택}이다($\Delta+$0.0136). 붙일 뻔했다. \textbf{대조군을
돌려서 살았다} --- 도서를 빼면 $+$0.0171, \textbf{펀딩을 빼면 $+$0.0191}로
애니보다 크다. 열 도메인 중 셋이 ``빼면 좋아진다''로 나오고, 그 크기가 그 도메인의
앙상블 $\rho$로 \textbf{$r=-$0.999} 예측된다. \textbf{산술이다} --- 판정치는
대상 평균이므로 평균보다 낮은 도메인은 무엇이든 빼면 오른다. 증거가 아니다.
그래서 \textbf{분모 불변 기준}을 새로 만들었다. 대상 집합을 $X$를 뺀 나머지로
\emph{고정}하고 $X$를 출처로 쓰느냐만 비교한다. 그렇게 재니 \textbf{열 도메인
전부 보류}다 --- $|\Delta|\le$0.0023이고 열 구간이 전부 0을 포함한다.
\textbf{어느 출처도 뺄 이유가 없다.} 애니(라프텔)도 그대로 둔다. 규약을 하나
더 둔다 --- \textbf{판정치는 설계 결정에는 쓰고 도메인 집합 결정에는 안 쓴다.}
\end{abstract}
]

\begin{figure*}[t]
\centering
\includegraphics[width=\textwidth]{figs/denom.pdf}
\caption{왼쪽 --- 도메인을 빼면 생기는 $\Delta$판정치. 오른쪽 --- 대상 집합을
고정한 출처 기여도.}
\end{figure*}

\section{붙일 뻔했다}

\begin{table}[t]
\centering\tsize
\caption{애니(라프텔) 제외. 씨앗 넷, 붓스트랩 150회.}
\begin{tabular}{lrl}
\toprule
씨앗 & $\Delta$ & 95\% 구간 \\
\midrule
20260729 & $+$0.0136 & $[+0.0028, +0.0194]$ \\
19770101 & $+$0.0136 & $[+0.0030, +0.0194]$ \\
20250315 & $+$0.0136 & $[+0.0023, +0.0204]$ \\
20260101 & $+$0.0136 & $[+0.0022, +0.0206]$ \\
\midrule
\multicolumn{2}{l}{세계애니 제외(대조)} & 0/4 \\
\bottomrule
\end{tabular}
\end{table}

씨앗 넷 전부 채택이다. 노트 72의 규약을 통과했고 노트 80에서 그 규약이 실제로
무언가를 막는 것도 봤다. 여기서 멈췄으면 애니를 뺐을 것이다.

\section{대조군이 살렸다}

\begin{table}[t]
\centering\tsize
\caption{열 도메인을 하나씩 빼 본다. 붓스트랩 100회.}
\begin{tabular}{lrl}
\toprule
& $\Delta$판정치 & 판정 \\
\midrule
\textbf{펀딩} 제외 & \textbf{$+$0.0191} & 채택 \\
\textbf{도서} 제외 & \textbf{$+$0.0171} & 채택 \\
\textbf{애니} 제외 & $+$0.0136 & 채택 \\
게임 제외 & $+$0.0001 & 보류 \\
세계애니 제외 & $-$0.0035 & 보류 \\
아이돌 제외 & $-$0.0085 & 보류 \\
만화 제외 & $-$0.0088 & 악화 \\
웹툰 제외 & $-$0.0094 & 악화 \\
모바일 제외 & $-$0.0242 & 악화 \\
\bottomrule
\end{tabular}
\end{table}

펀딩을 빼면 애니보다 더 오른다. 그리고 그 크기가 그 도메인의 앙상블 $\rho$로
\textbf{$r=-$0.999} 예측된다 --- 완전한 직선이다.

\claimx{판정치로 도메인 제거를 판정할 수 없다. 판정치는 대상 평균이므로 평균보다
낮은 도메인은 \emph{무엇이든} 빼면 오르고, 그 크기는 그 도메인의 성적으로
$r=-$0.999 예측된다. 산술이지 증거가 아니다.}{도메인을 빼면 셀도 줄어 구간이
넓어진다. 이 대조 실험은 평균만 봤고 정밀도 손실은 안 셌다 --- 그것까지 넣으면
셋 다 더 나쁘게 나올 것이다.}

\begin{quote}\fontsize{9.5}{11.5}\selectfont
\textbf{규약을 하나 더 둔다.} 판정치는 \emph{설계 결정}(배선 · 축 · 부분집합)에는
쓰고 \emph{도메인 집합 결정}에는 안 쓴다. 분모가 바뀌는 비교에는 다른 자가
필요하다.
\end{quote}

\section{분모 불변 기준}

대상 집합을 $X$를 뺀 나머지로 \textbf{고정}하고, 그 대상들을 맞힐 때 $X$를
출처로 쓰느냐만 비교한다. 분모가 같으므로 산술 효과가 없다.

\begin{table}[t]
\centering\tsize
\caption{출처 기여도. 대상 집합 고정, 붓스트랩 100회.}
\begin{tabular}{lrl}
\toprule
& $\Delta$ & 95\% 구간 \\
\midrule
팝업 & $+$0.0023 & $[-0.0011, +0.0061]$ \\
아이돌 & $+$0.0007 & $[-0.0024, +0.0051]$ \\
게임 & $-$0.0001 & $[-0.0028, +0.0027]$ \\
웹툰 & $-$0.0002 & $[-0.0037, +0.0044]$ \\
만화 & $-$0.0003 & $[-0.0034, +0.0027]$ \\
세계애니 & $-$0.0003 & $[-0.0031, +0.0028]$ \\
모바일 & $-$0.0006 & $[-0.0046, +0.0039]$ \\
도서 & $-$0.0010 & $[-0.0040, +0.0022]$ \\
애니 & $-$0.0012 & $[-0.0062, +0.0027]$ \\
펀딩 & $-$0.0022 & $[-0.0082, +0.0031]$ \\
\bottomrule
\end{tabular}
\end{table}

\textbf{열 도메인 전부 보류다.} 가장 큰 것이 $|\Delta|=$0.0023이고 열 구간이
전부 0을 포함한다.

\claimx{아홉 출처를 순위 평균하는 앙상블은 어느 한 출처를 빼도 흔들리지 않는다
--- $|\Delta|\le$0.0023. 나쁜 출처(애니 $-$0.0012)도 좋은 출처(팝업 $+$0.0023)도
마찬가지다. \textbf{뺄 것도 없고 뺄 필요도 없다.}}{출처가 아홉이라 하나의
비중이 9분의 1이다. 도메인이 넷이었을 때는 달랐을 것이고, 스물이 되면 더
안 흔들릴 것이다.}

이것이 노트 72의 희석 · 노트 73의 합의 정향 · 노트 80의 부분집합이 전부 같은
자리에서 막힌 이유이기도 하다. \textbf{앙상블은 한 출처에 관한 모든 처방을
9분의 1로 깎는다.}

\section{애니(라프텔)를 둔다}

노트 81이 물은 것에 답한다 --- \textbf{안 뺀다.} 팝업 출처로 $+$0.124인 것은
사실이지만, 대상 집합을 고정하고 재면 그것을 빼는 값이 $-$0.0012
$[-$0.0062, $+$0.0027$]$로 0과 구별되지 않는다. \textbf{나쁜 출처는 앙상블
안에서 이미 무해하다.}

\section{한계}

\begin{itemize}[leftmargin=1.1em,itemsep=0.15em]
\item 출처 기여도를 붓스트랩 100회로 쟀다. 노트 81이 반폭 비교는 300회로
      하라고 했는데 여기서는 안 지켰다 --- 다만 열 구간이 전부 0을 크게
      포함하므로 결론은 안 바뀔 것이다.
\item 한 번에 하나씩만 뺐다. 둘 이상 빼면 다를 수 있다.
\item ``도메인 제거는 판정치로 못 한다''를 보였는데, 그러면 \emph{도메인 추가}는
      어떤 자로 판정해야 하는가 --- 안 다뤘다. 지금까지 도메인 추가는 검정 없이
      해 왔다.
\item 출처 기여도가 전부 0이라는 것은 ``어느 출처도 필요 없다''로도 읽힌다.
      출처를 여러 개 동시에 빼면 무너질 텐데 그 경계를 안 찾았다.
\item 노트 81이 남긴 애니 매칭 확대는 진행 중이다(현재 83건, 32건에서 늘었다).
\end{itemize}

\section{다음}

\begin{enumerate}[leftmargin=1.3em,itemsep=0.15em]
\item \textbf{출처를 몇 개까지 빼도 되는가} --- 하나는 무해하다. 둘 · 셋 · 넷을
      빼며 경계를 찾는다. 그것이 ``도메인이 몇 개 필요한가''의 답이다.
\item 도메인 \emph{추가}를 판정할 자를 만든다. 분모가 바뀌는 것은 마찬가지다.
\item 애니 매칭 확대 후 라벨 신뢰도 재측정(노트 81의 32건 $\to$ 83건 이상).
\item 팝업 라벨의 신뢰도. 두 쌍이 0.72와 0.89로 갈리므로 팝업이 어느 쪽인지가
      제품 목표를 정한다.
\end{enumerate}

\vspace{0.6em}
\hrule height 0.3pt
\vspace{0.5em}
{\footnotesize
\textbf{참고문헌}\\[0.3em]
Simpson, E. H. (1951). The interpretation of interaction in contingency tables.
\emph{JRSS-B}, 13(2), 238--241.\\[0.15em]
Cawley, G. C., Talbot, N. L. C. (2010). On over-fitting in model selection and
subsequent selection bias in performance evaluation. \emph{JMLR}, 11, 2079--2107.\\[0.15em]
Breiman, L. (1996). Bagging predictors. \emph{Machine Learning}, 24(2), 123--140.\\[0.15em]
Kuncheva, L. I., Whitaker, C. J. (2003). Measures of diversity in classifier
ensembles. \emph{Machine Learning}, 51(2), 181--207.\\[0.15em]
Efron, B., Tibshirani, R. J. (1993). \emph{An Introduction to the Bootstrap}.
Chapman \& Hall.
}

\end{document}
